\documentclass{article}
\usepackage{iclr2027_conference,times}
\usepackage{amsmath,amssymb}
\usepackage{booktabs}
\usepackage{graphicx}
\usepackage{hyperref}
\usepackage{url}

\iclrfinalcopy

\title{Embedding Agentic Traces for Efficient Evaluation}

\author{Anonymous\thanks{Draft skeleton; author block TBD (Helivan, Jataware, Johns Hopkins).}\\
Helivan Research
}

\newcommand{\qubric}{\textsc{qubric}}

\begin{document}
\maketitle

\begin{abstract}
%% WORKING DRAFT ABSTRACT -- to be rewritten
Evaluating agentic systems is expensive: a single leaderboard entry
requires hundreds of long-horizon task executions, and modern systems are
sensitive to configuration details that multiply the number of variants
needing evaluation. We study the problem of predicting a new system's
full-benchmark score from a small number of probe runs, using publicly
cached runs of other systems as references. Each probe produces an
\emph{agentic trace} that carries far more information than its pass/fail
grade, but off-the-shelf semantic embeddings of traces are dominated by
evaluation-irrelevant confounders --- harness, model family, and exact
configuration --- rather than by the content of the work performed. We
introduce \qubric{}, a query-specific-rubric transformation that converts
off-the-shelf embedders into content-faithful, confounder-invariant trace
representations, and combine the resulting trace geometry with item
response theory in a simple blend. On SWE-bench Verified (107 systems,
100-task panel), the blend predicts full-benchmark resolve rates within
0.066 MAE from a single probe and 0.032 from twenty, is never worse than
the strongest correctness-only estimator, and is significantly better
exactly where graded evidence is scarce. The reference geometry is
constructed from public archives for approximately \$26.
\end{abstract}

\section{Introduction}
\label{sec:intro}

The dominant paradigm of generative model use is increasingly agentic:
systems that plan, invoke tools, observe results, and iterate over long
horizons, rather than emit a single response. Evaluation is a necessary
part of understanding an agent's capability and utility, and as agents
become more capable and are given more trust, the tasks on which they must
be evaluated become correspondingly complex. Recent agentic evaluations
--- SWE-bench and its variants \citep{jimenez2024swebench,
openai2024verified}, Terminal-Bench \citep{tbench2026}, ARC-AGI
\citep{chollet2019measure}, and multi-benchmark agent suites
\citep{kapoor2025hal} --- require up to hundreds of thousands of tokens
per task\footnote{In the SWE-bench Verified leaderboard archive we measure
a median of $\sim$40K tokens per trace, with a heavy tail exceeding 700K
tokens.} and evaluate each system on roughly one hundred to two thousand
tasks. At published per-task costs of \$0.13--\$4 (aggregate
\$1.84/rollout across 21{,}730 rollouts in \citealp{kapoor2025hal}), this
translates to roughly \$100--\$2{,}000 per full evaluation of a single
system, and tens of thousands of dollars for multi-benchmark sweeps. Given
the sensitivity of these systems to changes in system prompts, retrieval
databases, scaffold parameters, and model versions --- each of which
creates a new variant worth measuring --- these costs are untenable and
result in systematic \emph{under-evaluation}: most checkpoints, ablations,
and configurations are simply never scored \citep{kapoor2024agents}.

In other domains of machine learning evaluation, practitioners have
developed efficient techniques to reduce this resource burden: benchmark
compression and anchor-item selection \citep{vivek2024anchor,
maiapolo2024tinybenchmarks}, item response theory and adaptive testing
\citep{kipnis2025metabench, truong2025amortized}, and performance
prediction from observational records \citep{ruan2024observational,
zhuang2025embedllm}. One promising approach uses off-the-shelf semantic
embedding models to construct low-dimensional representations of the
systems themselves and casts benchmark prediction as an inference problem
in the low-dimensional space \citep{duderstadt2023data,
helm2025statistical, helm2026query}: it requires no labeled training data,
no white-box access, and exploits response caches that evaluation
ecosystems already produce. Unfortunately, off-the-shelf embedding
functions do not yet exist for agentic traces; nor are semantic embedding
functions, without modification, appropriate for embedding traces for the
purpose of efficient evaluation --- we show that they are dominated by
evaluation-irrelevant \emph{confounders} (harness, model family, exact
configuration) rather than by the content of the work performed. We
introduce \qubric{} (\emph{query-specific rubrics}) as a way to transform
these models into a useful tool for efficient system evaluation.

In this paper we make four contributions. First, we show that
off-the-shelf semantic embedding functions are not appropriate for
efficient evaluation of agentic systems: their trace representations
recover harness, model family, and even exact configuration at far above
chance while remaining near chance on outcome, and --- diagnostically ---
their prediction error is \emph{flat} in the number of probe traces.
Second, we show that extracting and embedding content from a trace
according to a query-specific rubric mitigates these issues: \qubric{}
collapses the confounder axes to approximately chance while retaining task
and behavioral fidelity, robustly across seven embedding models. Third, we
show that \qubric{} representations improve query-efficient evaluation:
blended with item-response-theoretic estimators, they are never worse than
the strongest correctness-only method at any probe budget and are
significantly better where graded evidence is scarce, predicting
full-benchmark resolve rates within 0.066 MAE from a single probe run.
Fourth, we release the reference geometry, code, and judged-trace corpus
--- constructed from public archives for approximately \$26 --- as
infrastructure for trace-based evaluation research.

\subsection{Problem statement}
\label{sec:problem}

We adapt the query-efficient evaluation framework of
\citet{helm2026query} to the agentic setting. Let $\mathcal{T} = \{t_1,
\ldots, t_N\}$ be a benchmark of $N$ tasks and let an agentic system $A$
be a (stochastic) map from tasks to \emph{traces}: executing $A$ on task
$t$ yields a trace $\tau = A(t)$, a long, semi-structured log of the
system's actions, observations, and outputs, together with a grade
$g(\tau) \in \{0, 1\}$ assigned by the benchmark harness. The quantity of
interest is the system's benchmark score
\[
y(A) \;=\; \frac{1}{N} \sum_{i=1}^{N} \mathbb{E}\left[ g(A(t_i)) \right],
\]
estimated exactly only by executing and grading $A$ on all $N$ tasks.

We are given a \emph{reference cache}: systems $A_1, \ldots, A_M$ that
have already been evaluated in full, so that for each reference $j$ we
observe its traces $\{\tau_{ij}\}_{i=1}^{N}$, its grades, and hence
$y(A_j)$. Such caches exist publicly at scale.\footnote{The SWE-bench
Verified leaderboard archive contains trajectories and outcomes for over
one hundred systems; we estimate it embodies \$50K--100K of agent compute
contributed by submitters.} A \emph{new} system $A_0$ is executed on only
a probe subset $S \subset \mathcal{T}$ with $m = |S| \ll N$, yielding
probe traces $\{\tau_{i0}\}_{i \in S}$ and (optionally) their grades. The
goal is an estimator $\hat{y}(A_0)$ minimizing $\mathbb{E}\,|\hat{y}(A_0)
- y(A_0)|$ at a given probe budget $m$, where the estimator may use the
probe traces, the probe grades, and the full reference cache.

Two aspects distinguish the agentic setting from the short-response
setting of \citet{helm2026query}. First, the observation associated with
each probe is not a short answer but a trace of $10^4$--$10^6$ tokens in
one of many mutually incompatible harness formats; whether and how this
observation can be embedded usefully is the central question of this
paper. Second, the reference population is confounded: references sharing
the target's underlying LLM or harness permit trivially accurate
prediction by score-copying, so we evaluate under a
\emph{leave-one-LLM-out} protocol --- references sharing the target's
underlying LLM are excluded --- and verify robustness under
leave-one-harness-out. The probe subset $S$ may be fixed, drawn at
random, or chosen per-target by adaptive item selection; we report all
three regimes.

\section{Related work}
\label{sec:related}

\paragraph{Agents, traces, and evaluations.}
Agentic benchmarks pair long-horizon tasks with automated harnesses:
SWE-bench and its verified subset \citep{jimenez2024swebench,
openai2024verified} for repository-level software issues, executed by
scaffolds such as SWE-agent \citep{yang2024sweagent} and OpenHands
\citep{wang2025openhands}; Terminal-Bench \citep{tbench2026} for
command-line work; and suites aggregating many such settings
\citep{kapoor2025hal}. The traces these systems emit have recently become
objects of study in their own right: symbolic action-alphabet analyses of
leaderboard trajectories \citep{mehtiyev2026beyond}, canonical action
taxonomies that document per-scaffold format fragmentation
\citep{shu2026traceprobe}, failure taxonomies \citep{cemri2025mast},
trajectory-level judging \citep{zhuge2025agentjudge}, and trajectory
compression \citep{xiao2025agentdiet}. This literature produces
taxonomies, statistics, and reviews rather than reusable vector
representations, and none of it addresses the confounder dominance that
its own attribution results imply: coding-agent trajectories are
attributable to scaffold and model at 85.7\% against an 11\% baseline
\citep{oderinwale2026trajectories}, and model-level fingerprints survive
paraphrase and summarization \citep{sun2025idiosyncrasies, wu2026dna}.
Benchmark-integrity findings --- memorization inflating resolve rates
\citep{liang2025illusion} and mislabeled gold patches \citep{yu2025utboost}
--- further motivate representations of \emph{how} systems work rather
than reliance on the pass/fail bit alone.

\paragraph{Efficient evaluations.}
Efficient evaluation methods exploit the approximate low-rank structure of
the model$\times$item success matrix: anchor items and benchmark
compression \citep{vivek2024anchor, maiapolo2024tinybenchmarks,
kipnis2025metabench}, amortized and adaptive item response theory
\citep{truong2025amortized}, and observational performance prediction
\citep{ruan2024observational, zhuang2025embedllm}. These consume only
correctness patterns --- a few bits per probe --- which imposes a
documented resolution floor at extreme budgets \citep{yauney2026reliable}.
Agent-specific efficiency work is nascent: subset selection collapses
under scaffold shift \citep{ndzomga2026efficient}, and psychometric models
extended with task features still consume correctness rather than trace
content \citep{ge2026psychometrics}. Closest to this work is the data
kernel perspective space lineage \citep{duderstadt2023data,
helm2025statistical, acharyya2024consistent}, and in particular
query-efficient evaluation from cached responses \citep{helm2026query},
which predicts a new model's score by comparing few probe responses
against cached reference coverage. We generalize that framework to the
agentic setting, where each probe observation is a trace rather than a
short response, and where the reference population's confounder structure
requires an invariant representation and a conservative reference
protocol.

\paragraph{Embeddings of non-standard objects.}
General-purpose text embedders \citep{reimers2019sbert} fix one notion of
similarity at training time; no single embedding can represent all
relevance relations \citep{weller2025limit}, and instruction-conditioned
variants \citep{su2023instructor} honor only coarse, trained-in task
tags. Conditional similarity \citep{deshpande2023csts} formalizes
similarity-under-a-criterion at sentence scale, and a family of
QA-mediated representations embeds LLM-generated text in place of the
original object: answers to an instruction \citep{peng2024inbedder},
answer vectors to question batteries \citep{benara2024qaemb}, and
description-based similarity \citep{ravfogel2024description}. Embeddings
of interaction data exist for short dialogues \citep{liu2022dial2vec} and
RL state-action trajectories \citep{ge2025trajectory} --- the latter
observing that trajectory embeddings encode agent ``style,'' precisely
the confounder our setting must remove. \qubric{} sits at the unclaimed
intersection: inference-time, training-free, rubric-specified
re-representation of very long, heterogeneous logs, with an explicit
invariance mechanism (per-task consensus centering) and an evaluation
protocol that measures --- rather than assumes --- confounder
invariance.

\section{qubric}
\label{sec:method}
%% TODO: port method section (Model 1 / Model 2 / centering; configuration;
%% pruning and full-trace inputs; cost). See notes/summary.tex.

\section{Experiments}
\label{sec:experiments}
%% TODO: port from notes/summary.tex + figures/q100_tables.md:
%%   - requirements evidence (radar, heatmap, flat raw curve, stability)
%%   - efficient evaluation results (Tables 1-2, savings curve)
%%   - protocol robustness (leave-one-harness-out)

\section{Discussion}
\label{sec:discussion}
%% TODO

\bibliography{paper_refs}
\bibliographystyle{iclr2027_conference}

\end{document}
